\documentclass[conference]{IEEEtran}
\IEEEoverridecommandlockouts
\usepackage{cite}
\usepackage{amsmath,amssymb,amsfonts}
\usepackage{algorithm}
\usepackage{algorithmic}
\usepackage{graphicx}
\usepackage{textcomp}
\usepackage{xcolor}
\usepackage{booktabs}
\usepackage{multirow}
\usepackage{url}

\begin{document}

\title{Visual Grounding for Automated Reward Design:\\
When Does VLM Feedback Help LLM-Generated Reward Functions?}

\author{\IEEEauthorblockN{Anonymous Authors}}

\maketitle

% =============================================================
\begin{abstract}
Designing reward functions for reinforcement learning remains a major bottleneck
in robot manipulation. Recent work (Eureka, Ma et al.\ 2023) shows that large
language models (LLMs) can generate and iteratively refine reward code from task
descriptions and scalar training statistics. We extend this paradigm by closing
the loop with a vision-language model (VLM) that observes rollout videos of the
best candidate policy and provides behavioral diagnosis to the LLM before each
reward revision.

We evaluate the VLM+LLM method on an 8-task ManiSkill suite spanning solved,
partially solved, unsolved, and ceiling regimes: PushCube, PickCube, PushT,
UnitreeG1PlaceAppleInBowl, OpenCabinetDoor, OpenCabinetDrawer,
AnymalC-Reach (3 VLM+LLM runs), and PegInsertionSide. Direct Eureka-style
LLM-only baselines are available for PushT, UnitreeG1, and PegInsertionSide;
the remaining tasks are used to characterize how VLM utility changes across
solution phases. Our main findings are:
(1)~On already-solved tasks (PushCube, PickCube), VLM feedback is unnecessary
and can even cause temporary regression by prompting over-engineered reward
designs.
(2)~On reward-actionable unsolved tasks (PushT, UnitreeG1), VLM feedback provides
high-value behavioral diagnosis---``press-and-freeze'' in PushT and
stage-specific bottlenecks in UnitreeG1---that scalar metrics alone cannot
reveal, leading to strong gains over Eureka (PushT: 37.5\% SE / 56.25\% SO vs
6.25\%; UnitreeG1: 87.5\% SE vs 12.5\%).
(3)~On partially solved tasks, VLM is mixed but still informative. Cabinet
tasks gain early and then plateau (Door: 75.0\%$\to$87.5\%; Drawer:
62.5\%$\to$75.0\%), while AnymalC improves sharply in early iterations but can
later regress through over-regularization (across 3 runs: Iter~0 SE
31.25--56.25\%, best run reaches 100\%).
(4)~PegInsertionSide exposes a diagnosis-actionability gap. VLM feedback
correctly tracks behavioral progress (reach$\to$hover$\to$grasp-and-hold) and
raises return (109.5$\to$221.8), yet both VLM+LLM and Eureka remain at 0\%
success because the binary grasp/alignment cliff is not bridged by reward
shaping within the current search budget.

These results suggest that VLM-in-the-loop reward design is governed not only
by the current \emph{solution phase} (solved / partially solved / unsolved) but
also by whether the diagnosed failure mode is \emph{reward-actionable}, i.e.,
whether it can be translated into a reward modification that PPO can actually
exploit under the available optimization budget.
\end{abstract}

\begin{IEEEkeywords}
reward design, vision-language models, reinforcement learning, manipulation,
LLM-guided optimization
\end{IEEEkeywords}

% =============================================================
\section{Introduction}

Reward engineering is one of the most labor-intensive aspects of applying
reinforcement learning (RL) to robotic manipulation. Hand-crafted reward
functions require domain expertise, extensive tuning, and often fail to
generalize across tasks. Recent advances in large language models (LLMs) have
opened a new avenue: using LLMs to \emph{automatically generate} reward
function code from natural-language task descriptions.

Eureka~\cite{ma2024eureka} demonstrated that an LLM can produce reward functions that
match or exceed human-designed ones across a range of tasks. The key insight is
an evolutionary outer loop: the LLM generates $K$ candidate reward functions per
iteration, each is trained with PPO in parallel, and the best candidate (by task
success rate) is selected. A ``reward reflection'' summary of training
statistics is fed back to the LLM for the next iteration.

However, a fundamental limitation of the Eureka approach is that the LLM
receives only \emph{scalar} feedback (success rate, return, per-component reward
statistics). When the task fails, the LLM knows \emph{that} it failed but not
\emph{how} or \emph{why}. For multi-stage manipulation tasks, this gap can be
critical: a scalar success rate of 0\% provides no signal about whether the
robot failed at reaching, grasping, transporting, or placing.

We propose closing this gap with a vision-language model (VLM) that observes
rollout videos of the best candidate's behavior and provides a structured
behavioral diagnosis. This diagnosis---identifying specific failure modes,
behavioral patterns, and suggested reward modifications---is injected into the
LLM's prompt alongside the scalar statistics before each reward revision.

Our contributions are:
\begin{enumerate}
  \item A VLM-in-the-loop extension of Eureka-style reward generation that adds
    visual behavioral feedback to the LLM's iterative refinement process
    (Section~\ref{sec:method}).
  \item An empirical evaluation across eight tasks and multiple solution
    regimes, revealing that VLM feedback is most valuable when failures are both
    unresolved and reward-actionable, and can be counterproductive on
    already-solved tasks
    (Section~\ref{sec:results}).
  \item Analysis of VLM feedback quality across solution phases,
    identifying failure modes of the VLM itself (severity miscalibration,
    repetitive diagnosis) and design guidelines for when to activate the visual
    feedback loop (Section~\ref{sec:analysis}).
\end{enumerate}


% =============================================================
\section{Related Work}

\subsection{LLM-based Reward Design}

Recent work has shown that LLMs can synthesize reward functions directly from
task descriptions. Language to Rewards~\cite{yu2023language} and
Text2Reward~\cite{xie2024text2reward} established this direction by mapping
language to reward parameters or executable reward code, but both still rely on
humans to interpret failures and revise the design.

Eureka~\cite{ma2024eureka} made the loop fully automatic by coupling reward-code
generation with evolutionary selection over PPO runs. Follow-up work extended
this recipe to sim-to-real transfer~\cite{ma2024dreureka}, tactile
manipulation~\cite{field2025text2touch}, and more sample-efficient objectives or
search procedures~\cite{sarukkai2024progresscounts,hazra2025revolve,gao2026rfagent}.
Other systems introduce stronger feedback structure, such as trajectory analysis
or separate generator-evaluator roles~\cite{li2024auto,sun2025card}.

The closest prior works to ours are RE-GoT and
ELEMENTAL~\cite{yao2025regot,chen2025elemental}, which also bring visual signals
into reward design. Our focus is different. Rather than proposing another
general reward-generation framework, we keep the Eureka outer loop fixed and ask
when visual feedback provides information that scalar training statistics do
not. This lets us isolate the value of VLM diagnosis itself.

\subsection{VLMs for Rewarding and Diagnosis}

Another line of work uses VLMs to produce reward signals from visual
observations. Some methods score observations directly with image-language
similarity~\cite{rocamonde2024vlmrm,sontakke2023roboclip}, while others query
VLM preferences and fit an explicit reward model from those
judgments~\cite{wang2024rlvlmf,ghosh2025prefvlm}. Representation-driven methods
such as LIV, GVL, VLAC, SARM, and
ReWiND~\cite{ma2023liv,ma2025gvl,zhai2025vlac,chen2026sarm,zhang2025rewind}
treat visual-language models as value estimators, critics, or stage-aware reward
models.

More closely related are works that use VLMs for success detection or failure
reasoning rather than direct reward computation. Du et
al.~\cite{du2023successdetectors} cast success detection as visual question
answering, and AHA~\cite{duan2025aha} trains a VLM to explain manipulation
failures from trajectory observations.

Our role for the VLM is narrower and more interpretable. We do not ask it to be
the reward function or the critic. Instead, the VLM produces a textual diagnosis
of policy behavior, and the LLM remains responsible for editing explicit reward
code. This keeps the reward itself programmable while still grounding reward
revision in visual evidence.

\subsection{Iterative Language Agents with Feedback}

Our method also fits a broader pattern in which language agents improve through
interaction with the environment. In robotics, SayCan, Inner Monologue, and
Code as Policies show that plans or executable code become more reliable when
language reasoning is grounded in affordances, perception, and execution
feedback~\cite{ahn2022saycan,huang2023inner,liang2023code}. In general-purpose
LLM agents, Voyager, Reflexion, Self-Refine, and
CRITIC~\cite{wang2023voyager,shinn2023reflexion,madaan2023selfrefine,gou2024critic}
use environmental signals, verbal critiques, or tools to repair behavior over
multiple iterations.

We instantiate the same idea for reward engineering. The key question in our
setting is not whether feedback matters, but which feedback modality is
informative enough to revise reward code. Our results show that scalar training
statistics are often sufficient on easy tasks, whereas visual feedback becomes
useful when failures must be localized in behavior, especially in harder or
multi-stage manipulation.


% =============================================================
\section{Method}\label{sec:method}

\subsection{Overview}

Our method extends the Eureka outer loop with a VLM feedback channel. The
algorithm proceeds in $N$ iterations. Each iteration: (1)~the LLM generates $K$
candidate reward functions, (2)~each candidate is trained with PPO in parallel,
(3)~the best candidate is selected by task success rate, (4)~the VLM observes
a rollout of the best candidate and produces a behavioral diagnosis, and (5)~the
full history (code, metrics, VLM diagnosis) is fed back to the LLM.

\subsection{Algorithm}

\begin{algorithm}[t]
\caption{VLM-in-the-Loop Reward Generation}
\label{alg:vlm_loop}
\begin{algorithmic}[1]
\REQUIRE Task description $\mathcal{T}$, state docs $\mathcal{D}$, iters $N$, candidates $K$, RL steps $T$, LLM $\mathcal{L}$, VLM $\mathcal{V}$
\ENSURE Optimized reward function $R^*$
\STATE $\mathcal{H} \leftarrow \emptyset$, $R_{\text{best}} \leftarrow \texttt{sparse}$
\FOR{$n = 1$ \TO $N$}
    \STATE Build prompt $p_n$ from $\{\mathcal{T}, \mathcal{D}, R_{\text{best}}, \mathcal{H}\}$
    \STATE $\mathcal{C} \leftarrow \{R_1, \ldots, R_K\} \cup \{R_{\text{best}}\}$
    \COMMENT{$K$ new candidates + elite carry-over}
    \FOR{$R_k \in \mathcal{C}$ \textbf{in parallel}}
        \STATE Train $\pi_k$ with PPO for $T$ steps using $R_k$
        \STATE Evaluate $\pi_k$: $m_k = (\text{success\_once}, \text{success\_at\_end}, \text{return})$
    \ENDFOR
    \STATE $k^* \leftarrow \arg\max_k m_k.\text{success\_at\_end}$
    \STATE Record rollout of $\pi_{k^*}$, extract frames $\mathcal{F}$
    \STATE $c_{\text{vlm}} \leftarrow \mathcal{V}(\mathcal{F}, m_{k^*})$
    \COMMENT{Behavioral diagnosis}
    \STATE $\mathcal{H} \leftarrow \mathcal{H} \cup \{(n, R_{k^*}, m_{k^*}, c_{\text{vlm}}, \{m_k\}_k)\}$
\ENDFOR
\RETURN $R_{k^*}$
\end{algorithmic}
\end{algorithm}

\subsection{VLM Behavioral Diagnosis}

The VLM receives sampled frames from the best candidate's rollout along with
the episode's success label and scalar metrics. It is prompted to produce a
structured report with three sections: (1)~what the robot is currently doing,
(2)~what is going wrong (specific failure modes), and (3)~actionable reward
signal adjustments. This structured format ensures the diagnosis is directly
usable by the LLM for reward code modification.

\subsection{Elite Carry-Over}

Following Eureka, we carry the previous best candidate as an ``elite'' into the
next iteration's candidate pool. This provides a safety net: even if all new
LLM-generated candidates perform poorly, the best-so-far is preserved. Our
analysis shows this mechanism is critical when VLM feedback leads the LLM to
over-engineer solutions on easy tasks.

\subsection{Comparison: VLM+LLM vs Eureka (LLM-only)}

The Eureka baseline uses the same outer loop but without VLM feedback. The LLM
receives only scalar metrics and reward reflection statistics. Both methods use
identical hyperparameters ($K=4$, $N=5$, same PPO configuration and seed).


% =============================================================
\section{Experimental Setup}

\subsection{Tasks}

We evaluate on eight ManiSkill~\cite{gu2023maniskill2} tasks:

\begin{table}[htbp]
\caption{Task characteristics}
\label{tab:tasks}
\centering
\begin{tabular}{lccl}
\toprule
Task & Stages & Robot & Difficulty \\
\midrule
PushCube-v1 & 1 & Panda & Easy \\
PickCube-v1 & 2 & Panda & Easy \\
PushT-v1 & 1 & Panda & Medium \\
OpenCabinetDoor-v1 & 1 & Panda & Medium \\
OpenCabinetDrawer-v1 & 1 & Panda & Medium \\
AnymalC-Reach-v1 & 1 & AnymalC & Medium \\
UnitreeG1PlaceApple & 6 & G1 humanoid & Hard \\
PegInsertionSide-v1 & 5 & Panda & Very Hard \\
\bottomrule
\end{tabular}
\end{table}

\subsection{Comparison Coverage}

Direct VLM+LLM vs Eureka comparisons are available for PushT-v1,
UnitreeG1PlaceAppleInBowl-v1, and PegInsertionSide-v1. PushCube-v1,
PickCube-v1, OpenCabinetDoor-v1, and OpenCabinetDrawer-v1 currently serve as
VLM+LLM runs for phase characterization. For AnymalC-Reach-v1, we analyze
three independent VLM+LLM runs to characterize variability in a partially
solved regime; no clean Eureka counterpart is available there.

\subsection{Hyperparameters}

All experiments use $K=4$ candidates per iteration, $N=5$ outer iterations, and
$N_{\text{eval}}=16$ evaluation episodes per candidate. Fitness is
$\text{success\_at\_end}$ (fraction of episodes where the task is completed at
the final timestep). The LLM is GPT-4o; the VLM is also GPT-4o with vision.
Seed is fixed at 9351 for all runs.

% TODO: PPO hyperparameters table


% =============================================================
\section{Results}\label{sec:results}

\subsection{Main Results}

\begin{table}[htbp]
\caption{Eight-task evaluation summary}
\label{tab:main}
\centering
\begin{tabular}{lcccc}
\toprule
Task & Iter 0 SE & Best VLM SE & Comparator & Outcome \\
\midrule
PushCube & 100.0 & 100.0 & --- & Solved at Iter~0 \\
PickCube & 100.0 & 100.0 & --- & Solved at Iter~0 \\
PushT & 0.0 & 37.5 & Eur. 6.25 & Actionable unsolved \\
UnitreeG1 & 0.0 & 87.5 & Eur. 12.5 & Actionable multi-stage \\
Door$^\dagger$ & 75.0 & 87.5 & --- & Early gain, then plateau \\
Drawer$^\dagger$ & 62.5 & 75.0 & --- & Early gain, then plateau \\
AnymalC$^a$ & 31.25--56.25 & 100.0 & 3 VLM runs & Early gain, then instability \\
PegInsertion$^b$ & 0.0 & 0.0 & Eur. 0.0 & Diagnosis without solve \\
\bottomrule
\end{tabular}
\end{table}

\noindent{\footnotesize $^\dagger$Cabinet tasks are VLM+LLM-only phase-probing runs; both improve at Iter~1 and then plateau.}
\noindent{\footnotesize $^a$AnymalC summarizes three VLM+LLM runs; the latest run peaks at 87.5\% SE, while the best run reaches 100\%.}
\noindent{\footnotesize $^b$PegInsertionSide has a paired Eureka baseline but both methods remain at 0\% success.}

Across the full suite, four regimes emerge. Easy tasks (PushCube, PickCube) are
already solved at Iter~0 and do not need VLM feedback. Actionable unsolved
tasks (PushT, UnitreeG1) show the clearest benefit over Eureka. Partially
solved tasks split into an articulated plateau regime (Door, Drawer) and a
last-mile unstable regime (AnymalC). PegInsertionSide forms a ceiling regime:
VLM changes diagnosis quality and return, but not success.

On paired tasks, VLM+LLM beats Eureka strongly when the failure mode is
reward-actionable (PushT: 37.5\% SE / 56.25\% SO vs 6.25\%; UnitreeG1:
87.5\% SE vs 12.5\%), but not when the task is dominated by a hard binary and
precision cliff (PegInsertionSide: both 0\%, though VLM+LLM reaches 221.8 best
return vs Eureka's 113.7).

\subsection{Per-Iteration Trajectories}

\textbf{UnitreeG1PlaceAppleInBowl} (hardest task): VLM+LLM progresses from 0\%
(Iter~0) to 100\% success\_once at Iter~2 and 87.5\% success\_at\_end at
Iter~3. Each VLM observation identifies a new stage-specific bottleneck:
``grasps but fails to transport'' (Iter~0), ``nudge-and-abandon'' regression
(Iter~1), ``near-bowl freeze'' (Iter~2--3). The Eureka baseline plateaus at
6.25--12.5\% SE.

\textbf{PushT-v1}: VLM+LLM reaches 37.5\% success\_at\_end and 56.25\%
success\_once at Iter~4 after VLM identifies a ``press-and-freeze'' failure
pattern. Eureka reaches only 6.25\% on both SE and SO.

\textbf{PushCube-v1}: Achieves 100\% SE at Iter~0. Subsequent VLM feedback
identifies ``overshoot / over-pushing'' every iteration, prompting complex
geometric strategies that often re-evaluate lower. The best observed SE remains
100\% at Iter~0; later selected elites drop to 81.25--93.75\%. The apparent
regression may not be purely VLM-induced: the same Iter~0
policy re-evaluated at only 43.75\% in a later iteration, indicating that the
initial 100\% was likely inflated by N=16 evaluation noise rather than
representing true performance.

\textbf{PickCube-v1}: Maintains 93.75--100\% SE across all iterations. VLM
correctly reports ``no clear failure'' in 3/5 iterations, overreacts once
(``near-failure / brittle strategy'' on a 100\% SE episode). The final best
candidate's anti-stall code originated from VLM's one useful diagnosis.

\textbf{OpenCabinetDoor / OpenCabinetDrawer}: Both tasks begin partially solved
and improve immediately at Iter~1 (Door: 75.0\%$\to$87.5\%; Drawer:
62.5\%$\to$75.0\%), then plateau. In Drawer, repeated candidate timeouts reduce
effective search width to $K=2$ throughout; in Door, candidate attrition lowers
effective $K$ in later iterations. VLM remains descriptively correct, but
improvement saturates once the residual diagnosis stops changing.

\textbf{AnymalC-Reach-v1} (3 runs): This task starts partially solved from
Iter~0 with a large SO-SE gap. Iter~0 SE ranges from 31.25\% to 56.25\%, while
SO is already 81.25--100\%. The best run reaches 100\% SE at Iter~2, but the
latest run follows 31.25\%$\to$87.5\%$\to$18.75\%$\to$75.0\%$\to$6.25\%,
showing that VLM-guided fixes can help early and then over-regularize terminal
behavior once the task is mostly solved.

\textbf{PegInsertionSide-v1}: Neither VLM+LLM nor Eureka achieves any success
across 5 iterations (SE=SO=0 throughout). However, VLM+LLM nearly doubles best
return (221.8 vs 113.7) and the VLM diagnosis evolves from
reach-and-hover near the box to grasp-and-hold local optima. This indicates
behavioral progress that scalar success metrics do not show, while also exposing
an expressivity/budget ceiling for reward shaping on hard binary grasp plus
precision alignment tasks.


% =============================================================
\section{Analysis}\label{sec:analysis}

\subsection{When Does VLM Feedback Help?}

\begin{table}[htbp]
\caption{VLM effectiveness by task and solution phase}
\label{tab:vlm_effectiveness}
\centering
\begin{tabular}{lccl}
\toprule
Task & Iter 0 SE & VLM Effect & Diagnosis Type \\
\midrule
PushCube & 100\% & Regression$^*$ & Stylistic (non-blocking) \\
PickCube & 100\% & Neutral & Mostly ``no failure'' \\
PushT & 0\% & Improvement & Blocking failure mode \\
UnitreeG1 & 0\% & Major gain & Stage bottleneck \\
OpenCabinetDoor$^\dagger$ & 75\% & Early gain, then plateau & Partial-open stall / contact drift \\
OpenCabinetDrawer$^\dagger$ & 62.5\% & Early gain, then plateau$^\ddagger$ & Partial-open retreat / disengagement \\
AnymalC$^a$ & 31--56\% & Mixed: early gain, later instability & Terminal hold / overshoot / posture \\
PegInsertion$^b$ & 0\% & Informative diagnosis, no solve & Binary grasp / alignment cliff \\
\bottomrule
\end{tabular}
\end{table}

\noindent{\footnotesize $^*$Observed regression is confounded with N=16 evaluation noise (same policy re-evaluated $\pm$56pp).}
\noindent{\footnotesize $^\dagger$Supplementary VLM+LLM-only analysis (no clean Eureka counterpart).}
\noindent{\footnotesize $^\ddagger$Drawer had effective $K=2$ due repeated candidate timeouts; Door also had candidate attrition in late iterations.}
\noindent{\footnotesize $^a$AnymalC summarizes three VLM+LLM runs with materially different post-breakthrough behavior.}
\noindent{\footnotesize $^b$PegInsertion has a paired Eureka baseline, but neither method reaches non-zero success.}

The value of VLM feedback is better explained by the policy's \emph{solution
phase} than by task label alone. We define progress as:
\begin{itemize}
  \item \textbf{Unsolved phase:} low SE and low SO (policy rarely reaches the goal).
  \item \textbf{Partially solved phase:} SO high but SE clearly lower (policy reaches but does not reliably hold/finish).
  \item \textbf{Solved phase:} high SE with stable re-evaluation across iterations.
\end{itemize}
Under this definition, VLM feedback is typically most useful in unsolved and
partially solved phases, and often less useful (or harmful) in the solved phase.
However, the full 8-task suite shows that phase alone is not sufficient. By
\emph{reward-actionable}, we mean that the VLM diagnosis can be mapped to a
concrete reward change and that PPO can actually learn from that change within
the available compute budget. PegInsertionSide is unsolved, yet VLM does not
improve success because the diagnosed failure is not reward-actionable under the
current setup. Thus, VLM value depends on two conditions jointly: the policy
must still be failing, and the failure must be addressable by reward shaping
within the available optimization horizon. The cabinet runs and AnymalC further
refine the partially solved regime: some tasks plateau after one useful
diagnosis, while others become unstable when repeated VLM-guided fixes
over-regularize the terminal behavior.

\subsection{VLM Failure Modes}

We identify several recurring failure modes and limits of the VLM feedback loop:

\textbf{Severity miscalibration.} The VLM is prompted to find ``failure modes''
and produces diagnoses with similar urgency regardless of success rate. On
PushCube (100\% SE), it reports ``overshoot / lack of termination'' with the
same alarm as on UnitreeG1 (0\% SE). This causes the LLM to treat non-blocking
stylistic issues as critical failures requiring architectural reward changes.

\textbf{Repetitive diagnosis on simple tasks.} PushCube's VLM produces the same
``overshoot'' observation across all 5 iterations. Unlike UnitreeG1 where VLM
discovers qualitatively new bottlenecks each iteration (grasp $\to$ transport
$\to$ place $\to$ release), simple tasks have only one observable behavior
pattern.

\textbf{Over-engineering induction.} VLM diagnoses on successful tasks prompt
the LLM to generate complex reward strategies (explicit geometric staging,
curriculum gates, behind-pose requirements) that perform worse than the simple
potential-based reward that already works. On PushCube, every iteration produced
at least one 0\% SE candidate inspired by VLM-suggested complexity.

\textbf{Search-width confound in plateau regimes.} In cabinet runs, multiple
candidates failed to complete training (timeouts/runtime failures), reducing
effective $K$ and lowering the chance of finding improvements in later
iterations. Thus, late-iteration plateau should be interpreted as a joint effect
of diagnosis saturation and reduced search width, not solely as a VLM failure.

\textbf{Over-regularization in last-mile tasks.} AnymalC shows that repeated
failure-focused feedback can become harmful after the task is mostly solved.
Across three runs, early VLM-guided edits improve SE substantially, but later
penalty accumulation can destabilize the terminal hold behavior and drive
regression (e.g., the latest run peaks at 87.5\% and ends at 6.25\%).

\textbf{Diagnosis-actionability gap.} PegInsertionSide demonstrates that a VLM
can be descriptively correct yet practically ineffective. The diagnosis evolves
meaningfully (hover near box $\to$ no insertion push $\to$ grasp-and-hold local
optimum), but the LLM cannot convert these observations into a reward function
that overcomes the binary grasp and precision alignment cliff within the current
training budget.

\subsection{Reward-Task Misalignment vs Reward Hacking}

A recurring pattern across tasks is high return with low success. We note that
this is better described as \emph{reward-task misalignment} rather than
``reward hacking'': the agent rationally optimizes the reward function, but the
reward fails to adequately capture the task objective. The agent is not
exploiting loopholes; the reward is simply misspecified. VLM feedback helps
detect and diagnose such misalignment on hard tasks by observing the specific
behavior the agent adopts under the misaligned reward.

\subsection{N=16 Evaluation Noise}

With only 16 evaluation episodes, success rate measurements have a granularity
of 6.25\% and substantial stochastic variance. We observe elite carry-over
candidates whose SE fluctuates by up to 56 percentage points across
re-evaluations of the identical policy (PushCube: 100\%$\to$43.75\%). This
noise affects both methods equally but complicates iteration-level comparisons.
% TODO: Consider increasing to N=64 or N=100 for final experiments.

\subsection{Role of Elite Carry-Over}

The elite mechanism acts as a critical safety net against VLM-induced regression
on easy tasks. Without it, the best PushCube solution (100\% SE at Iter~0) would
have been permanently lost when VLM-guided candidates underperformed. In cabinet
runs, elite carry-over preserved the Iter~1 breakthroughs (Door 87.5\%, Drawer
75\%) during post-breakthrough plateau. In AnymalC Run~A, elite carry-over also
preserved the Iter~2 100\% solution against later VLM-induced over-design. On
hard tasks, the elite preserves gains between breakthrough iterations.


% =============================================================
\section{Discussion and Limitations}

\textbf{When to use VLM feedback.}
Our results suggest phase-aware and actionability-aware activation of VLM
feedback. Here, \emph{actionability-aware} means distinguishing between
diagnoses that merely describe the failure and diagnoses that can be turned into
reward modifications from which PPO can actually learn. In reward-actionable
unsolved tasks (PushT, UnitreeG1), VLM diagnosis can substantially improve
intervention targeting. In partially solved tasks, VLM is often best used
early: cabinet tasks benefit once and then saturate, while AnymalC can regress
if failure-focused edits keep accumulating after a breakthrough. In solved
phases (PushCube, PickCube), VLM more often emphasizes non-blocking issues and
induces over-engineering. An adaptive system should therefore gate VLM usage by
phase, monitor whether diagnoses are still changing, and stop once SE is stable
or repeated feedback no longer yields gains.

\textbf{Reward-shaping ceiling.}
PegInsertionSide shows that informative diagnosis is not enough. Some failures
are bottlenecked by hard binary prerequisites or precision requirements that are
poorly served by reward shaping alone under a fixed training budget. In such
cases, curriculum learning, privileged intermediate labels, demonstrations, or
more training may be necessary before VLM feedback can translate into success.

\textbf{Single-seed evaluation.}
All experiments use a single seed (9351). While K=4 candidates per iteration
provide diversity in reward function design (the primary experimental variable),
the PPO training stochasticity is not independently sampled across seeds. We
argue this is acceptable because the reward function diversity---not seed
diversity---drives the key comparisons, and N=16 evaluation episodes provide
some statistical coverage. However, larger-scale multi-seed validation is
needed.

% TODO: Eureka baseline VLM leak bug. The current Eureka runs for PushT and
% UnitreeG1 have a stale VLM comment from Iter 0 leaked into the LLM prompt
% (from resuming a VLM+LLM run). The bug has been fixed in code but the
% baselines need to be re-run for a clean comparison. Current numbers should
% be interpreted as "fresh iterative VLM" vs "stale/frozen VLM" rather than
% a perfect VLM vs no-VLM ablation.

\textbf{VLM prompt design.}
The current VLM prompt always asks for ``failure analysis,'' which biases
it toward finding problems even on successful episodes. A severity-aware
prompt that conditions on the current success rate could reduce
over-engineering on easy tasks.

\textbf{Evaluation granularity.}
N=16 episodes is insufficient for fine-grained performance comparison.
The observed elite re-evaluation noise ($\pm$25 percentage points) limits
our ability to detect small improvements or confirm regressions.


% =============================================================
\section{Conclusion}

We presented a VLM-in-the-loop extension of Eureka-style automated reward
design. By closing the loop between policy behavior observation (VLM) and reward
code generation (LLM), we enable the system to diagnose and fix failure modes
that scalar training statistics alone cannot reveal.

Our main finding is that \textbf{VLM feedback value is phase-dependent but not
phase-determined}:
\begin{itemize}
  \item On hard, multi-stage tasks (UnitreeG1: 6 stages), VLM provides
    stage-specific bottleneck identification that drives 7$\times$ performance
    improvement over the LLM-only baseline (87.5\% vs 12.5\%).
  \item On medium tasks (PushT), VLM identifies behavioral failure modes
    invisible to scalar metrics, enabling a large gain over Eureka
    (37.5\% SE / 56.25\% SO vs 6.25\%).
  \item On partially solved articulated tasks (OpenCabinetDoor/Drawer), VLM
    provides early gains (+12.5pp at Iter~1) but then plateaus; repetitive
    diagnoses and reduced effective search width limit further improvement.
  \item On partially solved last-mile tasks (AnymalC), VLM can still produce
    sharp early gains, but repeated feedback may over-regularize terminal
    behavior and destabilize later iterations.
  \item On easy, already-solved tasks (PushCube, PickCube), VLM feedback is
    unnecessary and may cause regression through over-engineering induction,
    though observed regressions are confounded with N=16 evaluation noise.
  \item On very hard precision tasks (PegInsertionSide), VLM remains
    diagnostically informative but does not improve success; this exposes a
    diagnosis-actionability gap rather than a simple lack of visual signal.
\end{itemize}

These results point to a design principle for VLM-in-the-loop systems:
\textbf{visual feedback should be conditionally activated by both solution
phase and reward-actionability}. The VLM excels at answering ``how is the robot
failing?'', but this only improves performance when the answer can be converted
into an optimizable reward modification under the available budget.

% Future work: (1) VLM_full_RewardPlot extension combining rollout images with
% reward component plots for richer VLM diagnosis; (2) severity-aware VLM
% prompting; (3) multi-seed validation; (4) additional tasks (PegInsertion,
% OpenCabinetDoor).


% =============================================================
\bibliographystyle{IEEEtran}
\bibliography{refs}

\end{document}
